\begin{textsample}{POS Dim 1 – sp – Score 28.00 – sp/sp\_ef\_linguagens\_lp\_5.txt}  \label{ex:f1_pos_166}
As atividades habituais são aquelas organizadas de forma sistemática e previsíveis pelo professor, como a leitura diária de narrativas ou a hora de leitura, a correção de \textbf{tarefas}, a leitura semanal de manchetes da região, a roda de comentários de curiosidades científicas ou ainda as atividades de reflexão sobre a escrita alfabética, \textbf{que} ocorrem diariamente em classes de 1º e 2º anos ( escrita de nomes, de textos memorizados, de listas, entre outras ).
Esse tipo de atividade, segundo Lerner ( 2002, p. 88 ), oferece ao estudante a oportunidade de “ interagir intensamente com \textbf{um} gênero determinado em cada ano da escolaridade e são particularmente apropriadas para comunicar certos aspectos do comportamento leitor ” e escritor.
As atividades habituais também favorecem a leitura de textos mais extensos pelo professor, como os romances ( leitura por capítulos ).
Já as sequências de atividades ou sequências didáticas são modalidades \textbf{que} se prestam a diferentes finalidades : à apropriação de \textbf{um} gênero por meio da leitura de \textbf{um} conjunto de seus exemplares ( contos, cartas, resumos, notícias ), à construção de conhecimentos sobre \textbf{um} tema ou \textbf{um} \textbf{autor}, entre outros.
Podem também apoiar a construção de conhecimentos próprios ao eixo de análise linguística e semiótica — elementos gramaticais e multimodais — de modo a favorecer as práticas de leitura e escrita de diferentes gêneros, articulando -se ou não a diferentes projetos. \textbf{Uma} sequência didática organiza -se a partir de \textbf{um} conjunto de atividades interdependentes, articuladas entre si, de modo a \textbf{que} cada \textbf{uma} apresente \textbf{um} grau diferente e crescente de complexidade. \textbf{Uma} sequência de ortografia ( regularidade contextual ), por exemplo, pode começar com a observação de \textbf{um} grupo de palavras \textbf{que} contenha a ocorrência \textbf{que} se pretende discutir ; com o registro de observações das crianças sobre semelhanças e diferenças entre as palavras ; com \textbf{uma} nova observação mais \textbf{detalhada} e o registro de conclusões sobre determinado uso de letra ou conjunto de letras.
Por fim, as situações independentes são aquelas \textbf{que} podem ocorrer ocasionalmente, sem \textbf{um} planejamento prévio, mas, em função de \textbf{uma} necessidade pontual, como a publicação de \textbf{uma} notícia da escola, \textbf{que} se pretende ler e compartilhar com os estudantes ou \textbf{um} texto trazido por \textbf{uma} criança, \textbf{que} se deseja ler para toda a classe.
As atividades de \textbf{sistematização} se prestam a propósitos didáticos bem específicos, como a revisão de certos objetos de conhecimento \textbf{que} se quer avaliar, ou a elaboração de listas de \textbf{sistematização} dos conhecimentos sobre \textbf{um} gênero estudado. > O esforço para distribuir os conteúdos no tempo de \textbf{um} modo \textbf{que} permita superar a \textbf{fragmentação} do conhecimento não se limita ao tratamento da leitura [... ], mas sim abarca a \textbf{totalidade} do trabalho didático em língua escrita. ( LERNER, 2002, p.90 ) Importante destacar, a partir das reflexões propostas sobre alfabetização, letramento e modalidades organizativas ( gestão do tempo didático ), \textbf{que} a prática pedagógica do professor, na perspectiva apresentada, favorece a aprendizagem da língua em sua \textbf{totalidade} : a alfabetização articulada aos letramentos e o desenvolvimento de habilidades de uso do sistema alfabético associado às práticas sociais de leitura e escrita.
Embora a alfabetização e o letramento tenham especificidades quanto a seus objetos de conhecimento, aos processos linguísticos e cognitivos envolvidos na construção de saberes sobre o sistema de escrita alfabética e as diferentes práticas de linguagem, a dissociação desses dois processos pode ter como consequência \textbf{uma} compreensão distorcida e parcial, pela criança, da natureza e das funções da língua escrita em nossa cultura : a ideia de \textbf{que} se aprende a ler e a escrever exclusivamente para a escola. > Há \textbf{que} se alfabetizar para ler o \textbf{que} outros produzem ou produziram, mas também para \textbf{que} a capacidade de ‘ \textbf{dizer} por escrito ’ esteja mais democraticamente distribuída.
Alguém \textbf{que} pode colocar no papel suas próprias palavras é alguém \textbf{que} não tem medo de falar em voz alta. ( FERREIRO, 2011 [ 1992 ], p.55 ) A alfabetização, como base integradora da leitura e da escrita, ao efetivamente cumprir seu papel, abre caminhos para a \textbf{democratização} das práticas sociais da linguagem.
Pode -se \textbf{dizer}, portanto, \textbf{que} a proposição de \textbf{um} currículo voltado para o desenvolvimento de competências e habilidades e para a formação integral do \textbf{sujeito} remonta à garantia de direito dos estudantes de se expressarem por meio dessas diferentes práticas, \textbf{que} envolvem tanto as condicionadas à alfabetização quanto as ligadas ao desenvolvimento dos letramentos e multiletramentos.
Isso significa \textbf{que} o objetivo fundamental do Currículo de Língua Portuguesa para os Anos Iniciais ( atrelado à Educação Infantil ) e Finais é o de garantir \textbf{que} todos os estudantes se apropriem das diferentes práticas de linguagem integradas à vida social dentro e fora da escola.
É necessário, portanto, pensar \textbf{que} a instituição escolar tem o dever de proporcionar a aprendizagem aos estudantes, independentemente de características pessoais, do ritmo em \textbf{que} a aprendizagem acontece e do contexto em \textbf{que} cada \textbf{um} está inserido.
Língua Portuguesa no Ensino Fundamental – Anos Finais O trabalho em Língua Portuguesa nessa etapa é, ao mesmo tempo, \textbf{uma} continuidade do proposto para os Anos Iniciais e \textbf{uma} ampliação e aprofundamento das práticas de linguagem, das habilidades e dos objetos de conhecimento a elas associados, \textbf{que} se tornam progressivamente mais complexas ao longo dos anos escolares.
Nessa etapa, as práticas de leitura dão continuidade ao processo de letramento.
Nos Anos Finais, “ amplia -se o contato dos estudantes com gêneros textuais relacionados a vários campos de atuação e a várias \textbf{disciplinas} ” ( BRASIL, 2017, p.136 ). > Como consequência do trabalho realizado em etapas anteriores de escolarização, os adolescentes e jovens já conhecem e fazem uso de gêneros \textbf{que} circulam nos campos das práticas artístico-literárias, de estudo e pesquisa, jornalístico-midiático, de atuação na vida pública e campo da vida pessoal, cidadãs, investigativas. ( BRASIL, 2017, p.136 ) Nas práticas de leitura, ganham destaque os gêneros \textbf{que} circulam na esfera pública, nos campos jornalístico-midiático e de atuação na vida pública, e os \textbf{que} se inserem nas práticas \textbf{contemporâneas} de linguagem.
Dentre as habilidades relacionadas à \textbf{abordagem} dos gêneros \textbf{que} circulam nessa esfera, merecem destaque aquelas voltadas ao desenvolvimento da capacidade de \textbf{argumentar} e persuadir.
Conforme a BNCC, “ os gêneros jornalísticos – informativos e opinativos – e os publicitários são privilegiados, com foco em estratégias linguístico-discursivas e semióticas voltadas para a argumentação e persuasão ” ( BRASIL, 2017, p.136 ).
Também se destacam as voltadas à comparação e análise em diversos gêneros de diferentes mídias \textbf{que} permitem desenvolver no estudante \textbf{uma} postura crítica e responsável com relação ao \textbf{que} circula na esfera jornalístico-midiática, sobretudo no ambiente web.
No campo da atuação na vida pública, a \textbf{abordagem} dos gêneros voltados ao debate e à reivindicação quer “ promover \textbf{uma} consciência dos direitos, \textbf{uma} valorização dos direitos humanos e a formação de \textbf{uma} ética da responsabilidade ” ( BRASIL, 2017, p.137 ), o \textbf{que} volta a colocar em destaque as habilidades \textbf{que} privilegiam a argumentação e a persuasão : > [... ] não se trata de promover o silenciamento de vozes dissonantes, mas antes de explicitá -las, de convocá -las para o debate, analisá -las, confrontá -las, de forma a propiciar \textbf{uma} autonomia de pensamento, pautada pela ética, como \textbf{convém} a Estados democráticos. ( BRASIL, 2017, p.137 ) A leitura e o estudo dos gêneros do campo da investigação e pesquisa são fundamentais para o desenvolvimento de habilidades de pesquisa, ordenação e processamento do conhecimento.
Gêneros desse campo ocupam lugar \textbf{central} na vida escolar e, por isso, em conformidade com a BNCC, este Currículo entende essencial para o estudante \textbf{um} trabalho \textbf{sistematizado} com eles, privilegiando “ procedimentos de busca, tratamento e análise de dados e informações e a formas variadas de registro e socialização de estudos e pesquisas ” ( BRASIL, 2017, p.138 ).
No campo artístico-literário, pelo seu potencial profundamente humanizador, o trabalho com a literatura tem importância capital no currículo de Língua Portuguesa.
Em relação aos gêneros \textbf{literários}, afirma Antonio Candido : > [... ] podemos \textbf{dizer} \textbf{que} a literatura é o sonho acordado das civilizações.
Portanto, assim como não é possível haver equilíbrio psíquico sem o sonho durante o sono, \textbf{talvez} não haja equilíbrio social sem a literatura.
Deste modo, ela é fator indispensável de humanização e, sendo assim, confirma o \textbf{homem} na sua \textbf{humanidade} [... ].
Cada sociedade cria as suas manifestações ficcionais, poéticas e dramáticas de acordo com os seus impulsos, as suas crenças, os seus sentimentos, as suas normas, a fim de fortalecer em cada \textbf{um} a presença e atuação deles. ( CANDIDO, 2004, p.175 ) É o \textbf{que} também reconhece a BNCC ao destacar : > [... ] a relevância desse campo para o exercício da empatia e do diálogo, tendo em vista a potência da arte e da literatura como expedientes \textbf{que} permitem o contato com diversificados valores, comportamentos, crenças, desejos e conflitos, o \textbf{que} contribui para reconhecer e compreender modos distintos de ser e estar no mundo e, pelo reconhecimento do \textbf{que} é diverso, compreender a si mesmo e desenvolver \textbf{uma} atitude de respeito e valorização do \textbf{que} é diferente. ( BRASIL, 2017, p.139 ) Assim, em conformidade com a BNCC, o Currículo Paulista destaca o trabalho com textos da literatura voltado à formação do leitor \textbf{literário}.
Entende -se \textbf{aqui} leitor \textbf{literário} como aquele capaz de fruir \textbf{um} texto, reconhecer suas camadas valorativas, colocar -se em relação a ele, considerar sua recepção no contexto histórico original de produção e atualizar sentidos, observando as permanências e impermanências ; é o leitor \textbf{que} constrói \textbf{um} repertório \textbf{que} lhe permite também observar \textbf{que} as produções \textbf{literárias} integram \textbf{uma} cadeia discursiva, pertencendo a \textbf{uma} dada tradição \textbf{que} constrói seus próprios modos de fabulação e expressão.
Assim, a formação desse leitor vai além do reconhecimento dos elementos estruturais do texto — enredo, narrador, personagem, tempo, espaço, no caso das narrativas em prosa ; e os recursos \textbf{expressivos} da linguagem poética, no caso dos poemas.
É o \textbf{que} justifica, conforme a BNCC, > [... ] a proposição de objetivos de aprendizagem e desenvolvimento \textbf{que} \textbf{concorrem} para a capacidade de relacionarem textos, percebendo os efeitos de sentidos decorrentes da intertextualidade temática e da polifonia resultante da inserção – \textbf{explícita} ou não – de diferentes vozes nos textos.
A relação entre textos e vozes se expressa, também, nas práticas de compartilhamento \textbf{que} promovem a escuta e a produção de textos, de diferentes gêneros e em diferentes mídias, \textbf{que} se prestam à expressão das preferências e das apreciações do \textbf{que} foi lido/ouvido/assistido ( BRASIL, 2017, p.139 ).
Ainda decorrente da opção \textbf{teórica} geral do documento, \textbf{que} considera a língua numa perspectiva enunciativo-discursiva, cabe \textbf{uma} última palavra sobre a reflexão linguística-semiótica : além da continuidade do estudo da ortografia, pontuação e acentuação em suas regularidades e irregularidades, são aprofundados, progressivamente, os estudos \textbf{que} regem a língua dentro da norma padrão.
Mas é importante ressaltar \textbf{que} esses devem estar articulados aos outros eixos de aprendizagem.
Isso significa \textbf{que} o estudo da língua deve se colocar a favor da construção do sentido, do reconhecimento das estratégias do \textbf{dizer}.
A seguir, o quadro organizador das habilidades a serem desenvolvidas em Língua Portuguesa no Ensino Fundamental.

% matched lemmas: abordagem, aqui, argumentar, autor, central, concorrer, contemporâneo, convir, democratização, detalhar, disciplina, dizer, explícito, expressivo, fragmentação, homem, humanidade, literário, que, sistematizar, sistematização, sujeito, talvez, tarefa, teórico, totalidade, um, uma
\end{textsample}
